%==============================================================================
% Appendix B: Detailed Proof of the Pontryagin Protection Theorem
%==============================================================================
\section{Detailed Proof of the Pontryagin Protection Theorem}
\label{app:B}

This appendix derives the Pontryagin protection theorem of \S\ref{sec:pont}
algebraically.

\subsection{Statement}

\begin{theorem}[P0, restated]
  For the AX and VT modes and for the homogeneous minisuperspace topologies
  under consideration,
  $P = \innerprod{R}{\star R} \equiv 0$.
\end{theorem}

\subsection{AX Sector: Algebraic Proof of Pl\"ucker/Pfaffian Cancellation}
\label{app:B:AX}

\subsubsection*{Coframe definition}

The $\Sthree\times\Sone$ coframe used throughout this section is
\begin{equation}
  e^i = r\, M^i{}_j\,\sigma^j \quad (i=0,1,2),
  \qquad
  e^3 = L\!\left(d\tau + \omega_k\,\sigma^k\right),
\end{equation}
where $\sigma^i$ are the left-invariant 1-forms of $\Sthree$ and
$M^i{}_j = (\exp h)^i{}_j$ with $h = \sum_A q_A T_A$.  For isotropic AX mode
we set $h=0$ ($M=I$) and $\omega_k=\omega$ (all components equal).

The structure constants from the Maurer--Cartan equations
$de^k = (1/r)\varepsilon_{ijk}\,e^i\wedge e^j$,
$de^3 = (2L\omega/r^2)\varepsilon_{ij}\,e^i\wedge e^j$ are
\begin{equation}
  C^k_{ij} = \frac{2}{r}\varepsilon_{ijk},
  \qquad
  C^3_{ij} = \frac{4L\omega}{r^2}\varepsilon_{ij}.
\end{equation}
Defining
\begin{equation}
  c \equiv \frac{1+\eta}{r},
  \qquad
  f \equiv \frac{2L\omega}{r^2},
\end{equation}
the non-zero EC connection components are
\begin{equation}
  \omega^k_{ij} = c\,\varepsilon_{ijk},
  \qquad
  \omega^3_{ij} = f\,\varepsilon_{ij},
  \qquad
  \omega^i_{3j} = -f\,\varepsilon_{ij},
  \qquad
  \omega^i_{j3} = +f\,\varepsilon_{ij}.
\end{equation}

\subsubsection*{Riemann tensor formula for homogeneous spaces}

For homogeneous spaces the frame derivatives vanish and the Riemann tensor is
\begin{equation}
  R^a{}_{bcd}
  = \omega^a{}_{ec}\,\omega^e{}_{bd}
  - \omega^a{}_{ed}\,\omega^e{}_{bc}
  + \omega^a{}_{be}\,C^e{}_{cd}.
\end{equation}

\subsubsection*{Representative components for $(c,d)=(0,1)$}

The key Riemann components needed for the Pl\"ucker condition are
\begin{align}
  R_{0101} &= \frac{1-\eta^2}{r^2} + f^2, \\
  R_{2301} &= 0, \\
  R_{0201} &= -f^2, \\
  R_{1301} &= \frac{f(\eta-1)}{r}, \\
  R_{0301} &= \frac{f(1-\eta)}{r} = -R_{1301}, \\
  R_{1201} &= f^2.
\end{align}

\subsubsection*{Verification: Pfaffian $=0$ for $(c,d)=(0,1)$}

\begin{equation}
\begin{aligned}
  \mathrm{Pf}\,\omega^{(01)}
  &= R_{0101}\cdot R_{2301} - R_{0201}\cdot R_{1301} + R_{0301}\cdot R_{1201} \\
  &= \left(\tfrac{1-\eta^2}{r^2}+f^2\right)\cdot 0
   - (-f^2)\cdot\tfrac{f(\eta-1)}{r}
   + \tfrac{f(1-\eta)}{r}\cdot f^2 \\
  &= \frac{f^3(\eta-1)}{r} + \frac{f^3(1-\eta)}{r}
   = 0.
   \qquad\checkmark
\end{aligned}
\end{equation}

\subsubsection*{Verification for $(c,d)=(0,3)$ (S3-type)}

Using $C^e_{03}=0$,
\begin{equation}
  R_{0103}=-fc,\;
  R_{2303}=f^2,\;
  R_{0203}=-fc,\;
  R_{1303}=f^2,\;
  R_{0303}=-f^2,\;
  R_{1203}=0.
\end{equation}
\begin{equation}
  \mathrm{Pf}\,\omega^{(03)}
  = (-fc)\cdot f^2 - (-fc)\cdot f^2 + (-f^2)\cdot 0
  = -f^3c + f^3c = 0. \qquad\checkmark
\end{equation}

\subsubsection*{Generalisation to all $(c,d)$ pairs}

The six $(c,d)$ pairs fall into the SS-type $\{01,02,12\}$ and the S3-type
$\{03,13,23\}$.
\begin{itemize}
  \item \textbf{SS-type}: The cyclic symmetry of $\{0,1,2\}$ (the $SO(3)$ symmetry
        of $\Sthree$) gives the same algebraic structure as $(0,1)$.  The
        Pl\"ucker condition holds for the same algebraic reason.
  \item \textbf{S3-type}: As seen for $(0,3)$, the key is $C^e_{i3}=0$ (the
        $\tau$-direction generates no structure constants).  This holds for all
        S3-type pairs.
  \item $P_{33}$ \textbf{terms}: $R_{3\cdots 3\cdots}$ vanishes trivially by
        antisymmetry.
\end{itemize}
Hence the Pontryagin density is established algebraically:
\begin{equation}
  P = 2\sum_{(c,d)}\mathrm{Pf}\,\omega^{(cd)} = 0.
\end{equation}

\textbf{Remark.} The above proof is for isotropic twist.  For general vector
twist $(\omega_0,\omega_1,\omega_2)$ the same structure of structure constants is
preserved, and together with the identity $\Veff=f(I_2)$, it is confirmed
symbolically that $P\equiv 0$ (including $\Nilthree\times\Sone$).

\subsection{VT Sector: Pfaffian Cancellation}
\label{app:B:VT}

In the VT mode the cancellation mechanism differs.  For $\Sthree\times\Sone$,
\begin{equation}
  \sum_\mathrm{SS}\mathrm{Pf}\,\omega^{(cd)}
  = -\frac{16L^2 V\,I_2}{3r^5},
  \qquad
  \sum_\mathrm{S3}\mathrm{Pf}\,\omega^{(cd)}
  = +\frac{16L^2 V\,I_2}{3r^5},
\end{equation}
giving $P = \sum_\mathrm{SS}\mathrm{Pf} + \sum_\mathrm{S3}\mathrm{Pf} = 0$.

The cancellation is not trivial decomposability of each block (as in AX mode),
but an exact sign cancellation between the SS and S3 group Pfaffian totals.
The sign reversal originates from $V$ entering as pure vector-trace torsion only
in the $\tau$ direction, so the SS and S3 dual structures act complementarily.
The same cancellation holds algebraically for $\Nilthree\times\Sone$ as well.

\subsection{MX Sector: the Unique Exception}

In the MX mode $\eta\neq 0$ and $V\neq 0$ simultaneously, and the protection
breaks.  For $\Sthree\times\Sone$:
\begin{equation}
  P(\bm{\omega}) =
  \frac{2V\eta\,(V^2 r^2 + 9\eta^2 - 36)}{9r^3}
  + \frac{40V\eta L^2}{3r^5}\, I_2.
\end{equation}
Even at the isotropic point ($I_2=0$), $P_0 = 2V\eta(V^2 r^2+9\eta^2-36)/(9r^3)
\neq 0$, so the origin of $P\neq 0$ is torsion mixing $V\eta$, not the twist
$\omega$.

\subsection{NyVariant Independence}
\label{app:B:NyVar}

The Riemann tensor (and hence $P$) is independent of the NyVariant
(FULL/TT/REE) choice, as guaranteed theoretically by the pipeline computation
order.  The results by mode are:

\begin{table}[H]
\centering
\begin{tabular}{lccc}
\toprule
Mode & FULL & TT & REE \\
\midrule
AX & 0 & 0 & 0 \\
VT & 0 & 0 & 0 \\
MX & nonzero & nonzero & nonzero \\
\bottomrule
\end{tabular}
\caption{Pontryagin density by torsion mode and NyVariant.}
\label{tab:NyVar}
\end{table}

\subsection{Interpretation}

The key points of this appendix are the following.
\begin{enumerate}
  \item AX and VT are protected by distinct algebraic mechanisms: AX by
        individual cancellation per $(c,d)$ block (Pl\"ucker type); VT by sign
        cancellation between SS and S3 group Pfaffians.
  \item The MX mode is the unique exception; the twist itself is not the origin.
  \item The Pontryagin structure is cleanly separated sector by sector in a
        manner consistent with the final-basis organisation.
\end{enumerate}

\subsection{Computational Verification}

The algebraic computations of this appendix are confirmed symbolically by the
following scripts.
\begin{itemize}
  \item AX mode, Riemann block classification, per-$(c,d)$ Pfaffian zero check:\\
        \texttt{paper03/pontryagin\_anisotropic\_twist.py}
  \item VT/MX modes, NyVariant cross-check (all 18 cases):\\
        \texttt{paper03/torsion\_mode\_analysis.py}
  \item Isotropic twist $P\equiv 0$ scan ($r$, $\eta$ dependence):
        \texttt{paper03/pontryagin\_scan\_s3.py},
        \texttt{paper03/pontryagin\_scan\_nil3.py}
\end{itemize}
